% !TEX root = ../anexos.tex
\apendice{Especificación de Requisitos}

\section{Introducción}

Este apéndice recoge los requisitos del sistema, que se derivan de los objetivos del proyecto. Primero se enumeran los objetivos, después se recoge el catálogo de requisitos funcionales y no funcionales y, por último, se describen los casos de uso mediante tablas.

El usuario final es la única persona que interactúa con el sistema, a través de la aplicación web. Internamente, el sistema está formado por la aplicación web con un cliente que conecta con un modelo de lenguaje, y un servidor MCP que expone las herramientas de consulta y predicción del tráfico.

\section{Objetivos generales}

Este proyecto persigue construir una aplicación web que permita consultar, en lenguaje natural, el estado y la predicción del tráfico en distintas ciudades de España, apoyándose en modelos de redes neuronales de grafos (\textit{Graph Neural Networks}, GNN) y en el protocolo \textit{Model Context Protocol} (MCP) como capa de comunicación entre un modelo de lenguaje y los datos y modelos del sistema. A continuación se detallan los objetivos, agrupados por tipo.

\subsection{Objetivo general}

Desarrollar un sistema de apoyo a la movilidad urbana que responda a preguntas sobre el tráfico de ciudades españolas mediante una interfaz web, utilizando un modelo de predicción basado en GNN expuesto a través de un servidor MCP.

\subsection{Objetivos del software}

\begin{itemize}
  \item Permitir al usuario formular preguntas en lenguaje natural sobre el estado actual del tráfico en una ciudad y franja horaria concretas.
  \item Ofrecer predicciones de tráfico a corto plazo, indicando el horizonte temporal de la predicción.
  \item Presentar las respuestas de forma comprensible, indicando la ciudad, el momento y el nivel de congestión estimado.
  \item Soportar varias ciudades españolas, de modo que añadir una nueva no requiera modificar la lógica de la aplicación.
  \item Proteger los datos de la universidad: los datos originales permanecerán en el entorno donde se alojan.
\end{itemize}

\subsection{Objetivos técnicos}

\begin{itemize}
  \item Estudiar el estado del arte de MCP, de las redes LSTM y de las GNN aplicadas a la predicción de tráfico.
  \item Analizar, limpiar y preparar los datos de tráfico proporcionados por la universidad, definiendo la estructura de grafo (sensores y conexiones viarias) necesaria para las GNN.
  \item Implementar un modelo de referencia basado en LSTM y un modelo basado en GNN, y compararlos mediante métricas de error de predicción.
  \item Diseñar e implementar un servidor MCP en Python que exponga herramientas de consulta del estado del tráfico y de predicción, integrando los modelos entrenados.
  \item Desarrollar una aplicación web con un cliente MCP que conecte un modelo de lenguaje con el servidor, de forma que sea este quien decida qué herramienta invocar.
  \item Evaluar el sistema completo, tanto la calidad de las predicciones como el comportamiento de las respuestas a preguntas representativas de los usuarios.
\end{itemize}

\subsection{Objetivos de gestión y documentación}

\begin{itemize}
  \item Aplicar una metodología ágil con iteraciones (\textit{sprints}) de dos semanas, con objetivos revisados al final de cada una.
  \item Mantener el trabajo documentado y versionado en un repositorio de GitHub, con \textit{issues} y \textit{milestones} asociados a cada iteración.
  \item Elaborar la memoria del proyecto de forma continua durante el desarrollo, recogiendo las decisiones de diseño y los resultados.
\end{itemize}

\section{Catálogo de requisitos}

Los requisitos se identifican con el prefijo RF (funcionales) o RNF (no funcionales). La importancia puede ser alta, media o baja.

\subsection{Requisitos funcionales}

\tablaSmallSinColores{Requisitos funcionales}{l p{9cm} l}{rf}
{Id & Descripción & Importancia \\}
{
RF-1 & El sistema permitirá al usuario formular preguntas en lenguaje natural sobre el tráfico. & Alta \\
RF-2 & El sistema permitirá consultar las ciudades para las que hay datos de tráfico disponibles. & Media \\
RF-3 & El sistema permitirá consultar el estado del tráfico de una ciudad en una hora concreta, indicando la velocidad media y el nivel de congestión. & Alta \\
RF-4 & El sistema permitirá consultar la predicción del tráfico de una ciudad a corto plazo, indicando el horizonte de la predicción. & Alta \\
RF-5 & El sistema presentará las respuestas indicando la ciudad, el momento al que se refieren y el nivel de congestión estimado. & Alta \\
RF-6 & El sistema permitirá consultar información de la red viaria de una ciudad, como el número de sensores y los vecinos de un sensor. & Media \\
RF-7 & El sistema informará al usuario cuando falten datos o no sean válidos, y se los solicitará cuando proceda. & Media \\
RF-8 & El sistema mantendrá el contexto de la conversación para atender preguntas de seguimiento. & Media \\
}

\subsection{Requisitos no funcionales}

\tablaSmallSinColores{Requisitos no funcionales}{l p{9cm} l}{rnf}
{Id & Descripción & Importancia \\}
{
RNF-1 & Privacidad: los datos originales de la universidad permanecerán en su entorno. & Alta \\
RNF-2 & Fiabilidad: las respuestas se basarán exclusivamente en los resultados de las herramientas del servidor, sin inventar datos. & Alta \\
RNF-3 & Extensibilidad: añadir una nueva ciudad no requerirá modificar la lógica de la aplicación. & Media \\
RNF-4 & Usabilidad: las respuestas serán en español y comprensibles para usuarios sin conocimientos técnicos. & Media \\
RNF-5 & Mantenibilidad: el servidor MCP, los modelos y la aplicación web estarán separados, y el código estará versionado y documentado en GitHub. & Media \\
RNF-6 & Portabilidad: el sistema podrá ejecutarse en local con Python y las dependencias indicadas en el repositorio. & Baja \\
}

\section{Especificación de requisitos}

Los requisitos funcionales se concretan en los siguientes casos de uso, cuyo actor es el usuario de la aplicación web.

% Caso de Uso 1 -> Consultar las ciudades disponibles.
\begin{table}[p]
	\centering
	\begin{tabularx}{\linewidth}{ p{0.21\columnwidth} p{0.71\columnwidth} }
		\toprule
		\textbf{CU-1}    & \textbf{Consultar las ciudades disponibles}\\
		\toprule
		\textbf{Versión}              & 1.0    \\
		\textbf{Autor}                & Sergio Buil Laliena \\
		\textbf{Requisitos asociados} & RF-2 \\
		\textbf{Descripción}          & El usuario consulta para qué ciudades hay datos de tráfico disponibles. \\
		\textbf{Precondición}         & El servidor MCP está en ejecución y la aplicación puede conectarse a él. \\
		\textbf{Acciones}             &
		\begin{enumerate}
			\def\labelenumi{\arabic{enumi}.}
			\tightlist
			\item El usuario pide ver las ciudades disponibles.
			\item El modelo de lenguaje selecciona la herramienta de listado de ciudades.
			\item El servidor MCP devuelve la lista de ciudades.
			\item La aplicación muestra la lista al usuario.
		\end{enumerate}\\
		\textbf{Postcondición}        & El usuario conoce las ciudades sobre las que puede consultar. \\
		\textbf{Excepciones}          & Si el servidor no está disponible, se informa al usuario del error. \\
		\textbf{Importancia}          & Media \\
		\bottomrule
	\end{tabularx}
	\caption{CU-1 Consultar las ciudades disponibles.}
	\label{tabla:cu1}
\end{table}

% Caso de Uso 2 -> Consultar el estado del tráfico.
\begin{table}[p]
	\centering
	\begin{tabularx}{\linewidth}{ p{0.21\columnwidth} p{0.71\columnwidth} }
		\toprule
		\textbf{CU-2}    & \textbf{Consultar el estado del tráfico}\\
		\toprule
		\textbf{Versión}              & 1.0    \\
		\textbf{Autor}                & Sergio Buil Laliena \\
		\textbf{Requisitos asociados} & RF-1, RF-3, RF-5, RF-7 \\
		\textbf{Descripción}          & El usuario pregunta por el estado del tráfico de una ciudad en una hora concreta. \\
		\textbf{Precondición}         & El servidor MCP está en ejecución y la ciudad consultada dispone de datos. \\
		\textbf{Acciones}             &
		\begin{enumerate}
			\def\labelenumi{\arabic{enumi}.}
			\tightlist
			\item El usuario escribe una pregunta en lenguaje natural indicando la ciudad y la hora.
			\item El modelo de lenguaje interpreta la pregunta y selecciona la herramienta de estado del tráfico.
			\item El servidor MCP devuelve la velocidad media, el nivel de congestión y los sensores más lentos.
			\item El modelo de lenguaje redacta la respuesta a partir de esos resultados.
			\item La aplicación muestra la respuesta al usuario.
		\end{enumerate}\\
		\textbf{Postcondición}        & El usuario conoce el estado del tráfico de la ciudad a la hora consultada. \\
		\textbf{Excepciones}          & Si falta la ciudad o la hora, se las solicita al usuario. Si la ciudad no está disponible o la hora no es válida, se informa del error y, cuando procede, de las opciones válidas. \\
		\textbf{Importancia}          & Alta \\
		\bottomrule
	\end{tabularx}
	\caption{CU-2 Consultar el estado del tráfico.}
	\label{tabla:cu2}
\end{table}

% Caso de Uso 3 -> Consultar la predicción del tráfico.
\begin{table}[p]
	\centering
	\begin{tabularx}{\linewidth}{ p{0.21\columnwidth} p{0.71\columnwidth} }
		\toprule
		\textbf{CU-3}    & \textbf{Consultar la predicción del tráfico}\\
		\toprule
		\textbf{Versión}              & 1.0    \\
		\textbf{Autor}                & Sergio Buil Laliena \\
		\textbf{Requisitos asociados} & RF-1, RF-4, RF-5, RF-7 \\
		\textbf{Descripción}          & El usuario pregunta por la predicción del tráfico de una ciudad a corto plazo. \\
		\textbf{Precondición}         & El servidor MCP está en ejecución y la ciudad consultada dispone de un modelo de predicción. \\
		\textbf{Acciones}             &
		\begin{enumerate}
			\def\labelenumi{\arabic{enumi}.}
			\tightlist
			\item El usuario escribe una pregunta en lenguaje natural indicando la ciudad, la hora y el horizonte de la predicción.
			\item El modelo de lenguaje interpreta la pregunta y selecciona la herramienta de predicción.
			\item El servidor MCP ejecuta el modelo y devuelve el estado de tráfico previsto.
			\item El modelo de lenguaje redacta la respuesta indicando el momento al que se refiere la predicción.
			\item La aplicación muestra la respuesta al usuario.
		\end{enumerate}\\
		\textbf{Postcondición}        & El usuario conoce la previsión de tráfico de la ciudad para el momento consultado. \\
		\textbf{Excepciones}          & Si falta algún dato, se solicita al usuario. Si la ciudad no está disponible o el horizonte no es válido, se informa del error. \\
		\textbf{Importancia}          & Alta \\
		\bottomrule
	\end{tabularx}
	\caption{CU-3 Consultar la predicción del tráfico.}
	\label{tabla:cu3}
\end{table}

% Caso de Uso 4 -> Consultar la red viaria de una ciudad.
\begin{table}[p]
	\centering
	\begin{tabularx}{\linewidth}{ p{0.21\columnwidth} p{0.71\columnwidth} }
		\toprule
		\textbf{CU-4}    & \textbf{Consultar la red viaria de una ciudad}\\
		\toprule
		\textbf{Versión}              & 1.0    \\
		\textbf{Autor}                & Sergio Buil Laliena \\
		\textbf{Requisitos asociados} & RF-6 \\
		\textbf{Descripción}          & El usuario consulta información sobre la red de sensores de una ciudad, como su tamaño o los vecinos de un sensor. \\
		\textbf{Precondición}         & El servidor MCP está en ejecución y la ciudad consultada dispone de datos. \\
		\textbf{Acciones}             &
		\begin{enumerate}
			\def\labelenumi{\arabic{enumi}.}
			\tightlist
			\item El usuario pregunta por la red de sensores de una ciudad o por los vecinos de un sensor.
			\item El modelo de lenguaje selecciona la herramienta correspondiente.
			\item El servidor MCP devuelve el resumen del grafo o la lista de sensores vecinos.
			\item La aplicación muestra la respuesta al usuario.
		\end{enumerate}\\
		\textbf{Postcondición}        & El usuario conoce la estructura de la red viaria consultada. \\
		\textbf{Excepciones}          & Si la ciudad o el sensor no existen, se informa del error. \\
		\textbf{Importancia}          & Media \\
		\bottomrule
	\end{tabularx}
	\caption{CU-4 Consultar la red viaria de una ciudad.}
	\label{tabla:cu4}
\end{table}

% Caso de Uso 5 -> Realizar una consulta de seguimiento.
\begin{table}[p]
	\centering
	\begin{tabularx}{\linewidth}{ p{0.21\columnwidth} p{0.71\columnwidth} }
		\toprule
		\textbf{CU-5}    & \textbf{Realizar una consulta de seguimiento}\\
		\toprule
		\textbf{Versión}              & 1.0    \\
		\textbf{Autor}                & Sergio Buil Laliena \\
		\textbf{Requisitos asociados} & RF-1, RF-8 \\
		\textbf{Descripción}          & El usuario formula una pregunta que depende del contexto de una consulta anterior, por ejemplo cambiar la hora. \\
		\textbf{Precondición}         & Existe al menos una consulta previa en la conversación. \\
		\textbf{Acciones}             &
		\begin{enumerate}
			\def\labelenumi{\arabic{enumi}.}
			\tightlist
			\item El usuario escribe una pregunta que omite datos ya indicados antes.
			\item El modelo de lenguaje completa la ciudad y los demás datos con el contexto de la conversación.
			\item El modelo de lenguaje selecciona la herramienta adecuada y el servidor MCP devuelve el resultado.
			\item La aplicación muestra la respuesta al usuario.
		\end{enumerate}\\
		\textbf{Postcondición}        & El usuario obtiene la respuesta a su pregunta sin repetir los datos anteriores. \\
		\textbf{Excepciones}          & Si el contexto no es suficiente para completar la pregunta, se pide una aclaración al usuario. \\
		\textbf{Importancia}          & Media \\
		\bottomrule
	\end{tabularx}
	\caption{CU-5 Realizar una consulta de seguimiento.}
	\label{tabla:cu5}
\end{table}